%% SYNTHETIC TEST UNIT generated by writing/tools/make_outline_book.py from FINAL_OUTLINE_v2 -- not book text
\chapter{其他前沿严格结果}\label{ch:front}
本单元为测试用的合成文本（由 FINAL\_OUTLINE\_v2 自动生成），陈述内容从略。

\section{第 1 节（测试）}\label{sec:front-intro}

\begin{definition}[围道变形、钳夹、Landau 点与好方向]\label{def:front-trap}
（测试用合成陈述：内容从略。）
% uses: none
\end{definition}

\begin{lemma}[Gordan 二择一的构造性形式：最小范数方向]\label{lem:front-min-norm}
（测试用合成陈述：内容从略。）
% uses: none
\end{lemma}

\begin{remark}[Collins 的几何定理即 Gordan 二择一]\label{rem:front-gordan-collins}
（测试用合成陈述：内容从略。）
% uses: none
\end{remark}

\begin{example}[二维中的两种情形]\label{ex:front-gordan-2d}
（测试用合成陈述：内容从略。）
\uses{lem:front-min-norm}
\end{example}

\begin{proposition}["规避奇点需要严格正的一阶位移"不成立 \tagV]\label{prop:front-anomalous}
（测试用合成陈述：内容从略。）
\uses{def:front-trap}
\end{proposition}
\paragraph{地位与缺口}测试：见文献。

\section{第 2 节（测试）}\label{sec:front-landau}

\begin{theorem}[物理区钳夹判据（Collins 2020） \tagA]\label{thm:front-collins-pinch}
（测试用合成陈述：内容从略。）
\uses{def:front-trap, lem:appD-gordan, lem:front-min-norm, thm:ir-landau-momentum}
\end{theorem}
\begin{proofsketch}测试。\end{proofsketch}
\paragraph{地位与缺口}测试：见文献。

\begin{proposition}[Landau 条件只是奇点的必要条件 \tagV]\label{prop:front-landau-not-singular}
（测试用合成陈述：内容从略。）
\uses{def:front-trap}
\end{proposition}
\paragraph{地位与缺口}测试：见文献。

\begin{corollary}[动量空间中的 Coleman–Norton 图像（含无质量线与 Wilson 线） \tagR]\label{cor:front-coleman-norton}
（$R_{\mathrm{cond}}$：在所列编号假设之下证明完整；结论的地位不高于其中最弱的假设。）
前提：\cref{thm:front-collins-pinch}。
（测试用合成陈述：内容从略。）
\uses{thm:front-collins-pinch, prop:ir-coleman-norton}
\end{corollary}

\begin{example}[无质量自能的共线钳夹]\label{ex:front-collinear-pinch}
（测试用合成陈述：内容从略。）
\uses{cor:front-coleman-norton}
\end{example}

\begin{remark}[Coleman–Norton 与 ELOP 原证明的缺口]\label{rem:front-cn-history}
（测试用合成陈述：内容从略。）
% uses: none
\end{remark}

\section{第 3 节（测试）}\label{sec:front-ope}

\begin{proposition}[宽角振幅的红外森林公式（Ma 2020） \tagA]\label{prop:front-ir-forest}
（测试用合成陈述：内容从略。）
（本命题只涉及部分子态。）
\uses{def:ir-pinch-surface, prop:ir-power-counting, thm:ir-leading-regions, lem:method-ward-multi}
\end{proposition}
\begin{proofsketch}测试。\end{proofsketch}
\paragraph{地位与缺口}测试：本陈述的地位由\cref{thm:method-subtraction}、\cref{prop:sudakov-amplitude-factorization}决定。

\begin{proposition}[振幅的全阶硬—软—共线因子化：坐标空间路线与算符路线 \tagA]\label{prop:front-coordinate-space}
（测试用合成陈述：内容从略。）
（本命题只涉及部分子态。）
\uses{def:ir-pinch-surface, def:prelim-wilson}
\end{proposition}
\begin{proofsketch}测试。\end{proofsketch}
\paragraph{地位与缺口}测试：见文献。

\begin{proposition}[被积函数层次的局域红外因子化 \tagP]\label{prop:front-local-ir}
（测试用合成陈述：内容从略。）
\uses{thm:sudakov-ir-poles}
\end{proposition}
\paragraph{地位与缺口}测试：见文献。

\begin{remark}[代数几何的 Landau 分析与偶极子公式]\label{rem:front-landau-algebraic}
（测试用合成陈述：内容从略。）
% uses: none
\end{remark}

\begin{proposition}[标量场论中 OPE 的收敛、结合律与弯曲时空 \tagP]\label{prop:front-ope-scalar}
（测试用合成陈述：内容从略。）
\uses{thm:power-ope-convergence}
\end{proposition}
\paragraph{地位与缺口}测试：本陈述的地位由\cref{thm:power-ope-convergence}决定。

\section{第 4 节（测试）}\label{sec:front-regions}

\begin{proposition}[纯 Yang–Mills 理论中 OPE 在规范不变算符上闭合 \tagP]\label{prop:front-ope-ym}
（测试用合成陈述：内容从略。）
\uses{def:prelim-qcd, def:prelim-brst}
\end{proposition}
\paragraph{地位与缺口}测试：本陈述的地位由\cref{thm:dis-ope}决定。

\begin{proposition}[流方程（Wilson–Polchinski）重整化证明 \tagP]\label{prop:front-flow-renorm}
（测试用合成陈述：内容从略。）
\uses{def:prelim-qcd, def:prelim-brst}
\end{proposition}
\paragraph{地位与缺口}测试：见文献。

\begin{conjecture}[Yang–Mills 与 QCD 中 OPE 的收敛 \tagC]\label{conj:front-ope-convergence}
（测试用合成陈述：内容从略。）
\uses{thm:dis-ope, prop:front-ope-ym}
\end{conjecture}
\paragraph{地位与缺口}测试：见文献。

\begin{remark}[对含夸克的 QCD，本书能断言什么]\label{rem:front-ope-qcd-status}
（测试用合成陈述：内容从略。）
% uses: none
\end{remark}

\begin{lemma}[区域展开的主恒等式（Jantzen）]\label{lem:front-jantzen}
（测试用合成陈述：内容从略。）
\uses{def:ir-regions}
\end{lemma}

\section{第 5 节（测试）}\label{sec:front-glauber}

\begin{definition}[Lee–Pomeransky 表示、Newton 多面体与本质面]\label{def:front-facets}
（测试用合成陈述：内容从略。）
\uses{def:ir-landau-setup, def:ir-regions}
\end{definition}

\begin{proposition}[系数全正时的面展开（Semenova–Smirnov–Smirnov） \tagR]\label{prop:front-lp-facets}
（$R_{\mathrm{PT}}$：在重整化微扰论中对红外正规化的部分子态严格。）
（测试用合成陈述：内容从略。）
\uses{def:front-facets}
\end{proposition}

\begin{proposition}[宽角 Green 函数在壳展开：区域 = Landau 解 = 面 \tagR]\label{prop:front-onshell-facets}
（$R_{\mathrm{PT}}$：在重整化微扰论中对红外正规化的部分子态严格。）
（测试用合成陈述：内容从略。）
\uses{def:front-facets, def:ir-landau-setup, def:ir-pinch-surface}
\end{proposition}

\begin{assumption}[区域集合的完备性 \tagC]\label{asm:front-mor-complete}
（测试用合成陈述：内容从略。）
地位：测试状态句。
文献：测试。
被使用于：\cref{prop:front-commuting}。
解除途径：测试解除途径。
\uses{def:front-facets, def:ir-regions}
\end{assumption}

\begin{proposition}[多重在壳展开的可交换性 \tagR]\label{prop:front-commuting}
（$R_{\mathrm{cond}}$：在所列编号假设之下证明完整；结论的地位不高于其中最弱的假设。）
前提：\cref{asm:front-mor-complete}。
（测试用合成陈述：内容从略。）
\uses{asm:front-mor-complete, prop:front-onshell-facets}
\end{proposition}

\section{第 6 节（测试）}\label{sec:front-sll}

\begin{proposition}["区域 = Newton 多面体的下表面"在 Minkowski 运动学中不成立 \tagV]\label{prop:front-hidden-regions}
（测试用合成陈述：内容从略。）
\uses{def:front-facets}
\end{proposition}
\paragraph{地位与缺口}测试：本陈述的地位由\cref{rem:scet-assumptions}决定。

\begin{conjecture}[Minkowski 极限下区域方法的有效性 \tagC]\label{conj:front-mor-minkowski}
（测试用合成陈述：内容从略。）
\uses{def:front-facets, prop:front-lp-facets}
\end{conjecture}
\paragraph{地位与缺口}测试：本陈述的地位由\cref{thm:ir-regions-euclid}、\cref{prop:front-lp-facets}、\cref{lem:front-jantzen}、\cref{prop:ir-regions-minkowski}决定。

\begin{proposition}[Glauber 有效理论：Regge 化、BFKL 与零箱减除 \tagA]\label{prop:front-glauber-bfkl}
（测试用合成陈述：内容从略。）
（本命题只涉及部分子态。）
\uses{prop:scet-glauber-operators, thm:smallx-bfkl, def:scet-zero-bin}
\end{proposition}
\begin{proofsketch}测试。\end{proofsketch}
\paragraph{地位与缺口}测试：见文献。

\begin{proposition}[纯 Glauber 梯形图在单标度红外安全观测量中的相消 \tagA]\label{prop:front-syz}
（测试用合成陈述：内容从略。）
\uses{prop:front-glauber-bfkl, prop:scet-glauber-operators}
\end{proposition}
\begin{proofsketch}测试。\end{proofsketch}
\paragraph{地位与缺口}测试：见文献。

\begin{conjecture}[对多重部分子相互作用敏感的观测量：标准因子化的破坏 \tagC]\label{conj:front-mpi-violation}
（测试用合成陈述：内容从略。）
\uses{prop:front-syz}
\end{conjecture}
\paragraph{地位与缺口}测试：本陈述的地位由\cref{prop:front-syz}、\cref{prop:break-jet-veto}决定。

\section{第 7 节（测试）}\label{sec:front-nlp}

\begin{proposition}[重末态与旁观者之间的软末态相互作用 \tagA]\label{prop:front-heavy-fsi}
（测试用合成陈述：内容从略。）
\uses{prop:method-fsi}
\end{proposition}
\begin{proofsketch}测试。\end{proofsketch}
\paragraph{地位与缺口}测试：见文献。

\begin{remark}[Glauber 现状与本书分级的后果]\label{rem:front-glauber-status}
（测试用合成陈述：内容从略。）
% uses: none
\end{remark}

\begin{proposition}[非全局 LHC 观测量的因子化与超领头对数的全阶重求和 \tagA]\label{prop:front-sll-all-orders}
（测试用合成陈述：内容从略。）
（本命题只涉及部分子态。）
\uses{thm:jets-ngl-factorization, prop:sudakov-coulomb}
\end{proposition}
\begin{proofsketch}测试。\end{proofsketch}
\paragraph{地位与缺口}测试：本陈述的地位由\cref{prop:front-restoration}、\cref{prop:jets-sll}决定。

\begin{proposition}[超领头对数的塔与 $q\bar q$ 单态交换的闭式 \tagA]\label{prop:front-sll-closed-form}
（测试用合成陈述：内容从略。）
（本命题只涉及部分子态。）
\uses{prop:front-sll-all-orders}
\end{proposition}
\begin{proofsketch}测试。\end{proofsketch}
\paragraph{地位与缺口}测试：本陈述的地位由\cref{prop:front-sll-all-orders}决定。

\begin{example}[由闭式读出第一个超领头对数]\label{ex:front-sll-first}
（测试用合成陈述：内容从略。）
\uses{prop:front-sll-closed-form}
\end{example}

\section{第 8 节（测试）}\label{sec:front-nlp-tmd}

\begin{proposition}[Glauber 胶子引起的"因子化恢复" \tagP]\label{prop:front-restoration}
（测试用合成陈述：内容从略。）
（本命题只涉及部分子态。）
\uses{prop:front-sll-all-orders, thm:pdf-dglap}
\end{proposition}
\paragraph{地位与缺口}测试：本陈述的地位由\cref{prop:front-hidden-regions}决定。

\begin{proposition}[1-jettiness 中的相干性破坏对数 \tagV]\label{prop:front-jettiness-cvl}
（测试用合成陈述：内容从略。）
\uses{def:jets-algorithms, prop:sudakov-coulomb}
\end{proposition}
\paragraph{地位与缺口}测试：见文献。

\begin{assumption}[SCET 的模式完备性与 Glauber 模式的缺席 \tagA]\label{asm:front-scet-modes}
（测试用合成陈述：内容从略。）
地位：测试状态句。
文献：测试。
被使用于：\cref{thm:front-hgg}、\cref{prop:front-ggh}、\cref{prop:front-nlp-ll}、\cref{prop:front-nlp-dis}、\cref{prop:front-sidis-nlp}。
解除途径：测试解除途径。
\uses{def:scet-modes}
\end{assumption}
\begin{proofsketch}测试。\end{proofsketch}

\begin{definition}[端点发散与重因子化条件]\label{def:front-endpoint}
（测试用合成陈述：内容从略。）
\uses{def:scet-modes}
\end{definition}

\begin{theorem}[$h\to\gamma\gamma$（$b$ 夸克圈）的重整化端点因子化定理 \tagA]\label{thm:front-hgg}
（测试用合成陈述：内容从略。）
（本命题只涉及部分子态。）
\uses{asm:front-scet-modes, def:front-endpoint, prop:scet-scet2}
\end{theorem}
\begin{proofsketch}测试。\end{proofsketch}
\paragraph{地位与缺口}测试：本陈述的地位由\cref{asm:front-scet-modes}决定。

\section{第 9 节（测试）}\label{sec:front-positivity}

\begin{example}[领头双对数与 ${}_2F_2$]\label{ex:front-hgg-ldl}
（测试用合成陈述：内容从略。）
\uses{thm:front-hgg}
\end{example}

\begin{proposition}[轻夸克诱导的 $gg\to h$ 与重因子化减除方案 \tagA]\label{prop:front-ggh}
（测试用合成陈述：内容从略。）
（本命题只涉及部分子态。）
\uses{thm:front-hgg, asm:front-scet-modes}
\end{proposition}
\begin{proofsketch}测试。\end{proofsketch}
\paragraph{地位与缺口}测试：见文献。

\begin{proposition}[次领头幂次阈值 Drell–Yan 的 $d$ 维因子化 \tagP]\label{prop:front-nlp-threshold}
（测试用合成陈述：内容从略。）
（本命题只涉及部分子态。）
\uses{def:thr-limit, thm:thr-refactorization, def:front-endpoint}
\end{proposition}
\paragraph{地位与缺口}测试：本陈述的地位由\cref{prop:thr-nlp}决定。

\begin{proposition}[次领头幂次的领头对数重求和与端点因子化条件 \tagA]\label{prop:front-nlp-ll}
（测试用合成陈述：内容从略。）
（本命题只涉及部分子态。）
\uses{asm:front-scet-modes, def:front-endpoint, thm:thr-refactorization}
\end{proposition}
\begin{proofsketch}测试。\end{proofsketch}
\paragraph{地位与缺口}测试：本陈述的地位由\cref{conj:front-nlp-beyond-ll}决定。

\begin{proposition}[次领头幂次 DIS 与 PDF 因子化：阈值 PDF 的双标度结构 \tagA]\label{prop:front-nlp-dis}
（测试用合成陈述：内容从略。）
（本命题只涉及部分子态。）
\uses{asm:front-scet-modes, def:front-endpoint, def:pdf-quark, thm:pdf-dis-xspace}
\end{proposition}
\begin{proofsketch}测试。\end{proofsketch}
\paragraph{地位与缺口}测试：见文献。

\section{第 10 节（测试）}\label{sec:front-brst}

\begin{conjecture}[超出领头对数的次领头幂次端点因子化 \tagC]\label{conj:front-nlp-beyond-ll}
（测试用合成陈述：内容从略。）
\uses{def:front-endpoint, prop:front-nlp-ll}
\end{conjecture}
\paragraph{地位与缺口}测试：见文献。

\begin{conjecture}[次领头幂次 TMD 过程中 Glauber 交换的相消 \tagC]\label{conj:front-glauber-nlp}
（测试用合成陈述：内容从略。）
\uses{def:tmd-kin, thm:tmd-factorization}
\end{conjecture}
\paragraph{地位与缺口}测试：本陈述的地位由\cref{prop:front-tmd-nlp-operator}、\cref{prop:front-sidis-nlp}决定。

\begin{proposition}[TMD 算符展开与运动学幂次修正 \tagA]\label{prop:front-tmd-nlp-operator}
（测试用合成陈述：内容从略。）
（本命题只涉及部分子态。）
\uses{thm:tmd-factorization, def:tmd-collins}
\end{proposition}
\begin{proofsketch}测试。\end{proofsketch}
\paragraph{地位与缺口}测试：本陈述的地位由\cref{conj:front-glauber-nlp}决定。

\begin{proposition}[SIDIS 方位角不对称的次领头幂次 SCET 因子化 \tagA]\label{prop:front-sidis-nlp}
（测试用合成陈述：内容从略。）
\uses{thm:tmd-factorization, def:tmd-soft, def:scet-modes, asm:front-scet-modes}
\end{proposition}
\begin{proofsketch}测试。\end{proofsketch}
\paragraph{地位与缺口}测试：本陈述的地位由\cref{conj:front-glauber-nlp}、\cref{asm:front-scet-modes}决定。

\begin{proposition}[单圈次领头幂次系数函数与扭度 3 TMD 演化 \tagP]\label{prop:front-nlp-tmd-oneloop}
（测试用合成陈述：内容从略。）
（本命题只涉及部分子态。）
\uses{prop:front-tmd-nlp-operator}
\end{proposition}
\paragraph{地位与缺口}测试：见文献。

\section{第 11 节（测试）}\label{sec:front-smallx}

\begin{remark}[次领头幂次的幂律快度发散]\label{rem:front-power-rapidity}
（测试用合成陈述：内容从略。）
% uses: none
\end{remark}

\begin{proposition}[$\overline{\rm MS}$ 部分子分布可以为负 \tagV]\label{prop:front-msbar-negative}
（测试用合成陈述：内容从略。）
\uses{def:pdf-quark, def:prelim-msbar}
\end{proposition}
\paragraph{地位与缺口}测试：本陈述的地位由\cref{conj:pdf-positivity}决定。

\begin{example}[Yukawa 模型分布的端点行为与符号变化]\label{ex:front-yukawa}
（测试用合成陈述：内容从略。）
\uses{prop:front-msbar-negative}
\end{example}

\begin{proposition}[足够大标度上的正定性 \tagP]\label{prop:front-positivity-large-scale}
（测试用合成陈述：内容从略。）
非微扰地位：\cref{asm:intro-wightman}、\cref{asm:intro-composite}、\cref{asm:intro-bridge}。
\uses{thm:pdf-dis-xspace, rem:pdf-schemes}
\end{proposition}
\paragraph{地位与缺口}测试：见文献。

\begin{remark}[对本书与拟合的后果]\label{rem:front-positivity-tracks}
（测试用合成陈述：内容从略。）
% uses: none
\end{remark}

\section{第 12 节（测试）}\label{sec:front-dps}

\begin{remark}[为什么需要 BRST：K 胶子与非定域缺口]\label{rem:front-k-gluon}
（测试用合成陈述：内容从略。）
% uses: none
\end{remark}

\begin{proposition}[双局域 PDF 算符 BRST 分类的矩约化 \tagA]\label{prop:front-brst-bilocal}
（测试用合成陈述：内容从略。）
（本命题只涉及部分子态。）
\uses{def:pdf-lightray, prop:appJ-local-operators, thm:pdf-support-pt, prop:appL-bare, prop:appL-renormalized, thm:appK-muntz}
\end{proposition}
\begin{proofsketch}测试。\end{proofsketch}
\paragraph{地位与缺口}测试：本陈述的地位由\cref{prop:appL-renormalized}、\cref{conj:front-brst-multilocal}决定。

\begin{conjecture}[多局域光线算符的 BRST 分类 \tagC]\label{conj:front-brst-multilocal}
（测试用合成陈述：内容从略。）
\uses{def:prelim-brst, prop:appJ-local-operators}
\end{conjecture}
\paragraph{地位与缺口}测试：本陈述的地位由\cref{rem:front-k-gluon}、\cref{prop:lattice-renormalizability}、\cref{asm:lamet-aux-renorm}、\cref{rem:appJ-nonlocal}、\cref{prop:front-brst-bilocal}决定。

\begin{proposition}[NLO BK/JIMWLK 与三圈 BK 核 \tagA]\label{prop:front-nlo-bk}
（测试用合成陈述：内容从略。）
\uses{thm:smallx-bk, thm:smallx-bfkl}
\end{proposition}
\begin{proofsketch}测试。\end{proofsketch}
\paragraph{地位与缺口}测试：见文献。

\begin{proposition}[NLO BK 的不稳定性与共线重求和 \tagP]\label{prop:front-bk-collinear}
（测试用合成陈述：内容从略。）
\uses{prop:front-nlo-bk}
\end{proposition}
\paragraph{地位与缺口}测试：本陈述的地位由\cref{rem:front-open}决定。

\section{第 13 节（测试）}\label{sec:front-jets}

\begin{proposition}[高能极限中的 Wilson 线方法、胶子 Regge 化与 Regge 极限的红外奇异部分 \tagA]\label{prop:front-regge}
（测试用合成陈述：内容从略。）
（本命题只涉及部分子态。）
\uses{thm:smallx-bfkl, thm:sudakov-ir-poles}
\end{proposition}
\begin{proofsketch}测试。\end{proofsketch}
\paragraph{地位与缺口}测试：见文献。

\begin{proposition}[简单 Regge 极点因子化在 NNLL 失效 \tagV]\label{prop:front-regge-pole}
（测试用合成陈述：内容从略。）
\uses{def:smallx-regge, thm:sudakov-ir-poles}
\end{proposition}
\paragraph{地位与缺口}测试：本陈述的地位由\cref{thm:sudakov-ir-poles}、\cref{prop:front-regge}决定。

\begin{proposition}[小 $x$ 的 TMD 与 Sudakov 对数 \tagA]\label{prop:front-smallx-tmd}
（测试用合成陈述：内容从略。）
\uses{thm:smallx-bk}
\end{proposition}
\begin{proofsketch}测试。\end{proofsketch}
\paragraph{地位与缺口}测试：见文献。

\begin{proposition}[小 $x$ 螺旋度演化：截距与 BER 的分歧 \tagA]\label{prop:front-smallx-helicity}
（测试用合成陈述：内容从略。）
\uses{thm:smallx-bk, def:pdf-polarized}
\end{proposition}
\begin{proofsketch}测试。\end{proofsketch}
\paragraph{地位与缺口}测试：见文献。

\begin{proposition}[无重复计数的双部分子散射 \tagA]\label{prop:front-dps-scheme}
（测试用合成陈述：内容从略。）
非微扰地位：\cref{asm:intro-composite}、\cref{asm:intro-bridge}。
\uses{def:break-dpdf, thm:break-double-dy}
\end{proposition}
\begin{proofsketch}测试。\end{proofsketch}
\paragraph{地位与缺口}测试：本陈述的地位由\cref{conj:break-dps-coloured}决定。

\section{第 14 节（测试）}\label{sec:front-quarkonium}

\begin{proposition}[多部分子软因子快度发散的重整化定理 \tagA]\label{prop:front-rapidity-renorm}
（测试用合成陈述：内容从略。）
\uses{def:tmd-soft, def:break-dpdf}
\end{proposition}
\begin{proofsketch}测试。\end{proofsketch}
\paragraph{地位与缺口}测试：本陈述的地位由\cref{prop:yzliu-rapidity-linear}决定。

\begin{proposition}[色关联 DPD 的演化与匹配 \tagP]\label{prop:front-dpd-evolution}
（测试用合成陈述：内容从略。）
\uses{def:break-dpdf, prop:front-rapidity-renorm}
\end{proposition}
\paragraph{地位与缺口}测试：见文献。

\begin{proposition}[非全局对数：超出领头对数与测量驱动的因子化 \tagA]\label{prop:front-ngl}
（测试用合成陈述：内容从略。）
（本命题只涉及部分子态。）
\uses{thm:jets-ngl-factorization, thm:jets-bms}
\end{proposition}
\begin{proofsketch}测试。\end{proofsketch}
\paragraph{地位与缺口}测试：本陈述的地位由\cref{prop:jets-ngl-beyond}决定。

\begin{proposition}[喷注子结构、修饰喷注与半单举喷注函数 \tagA]\label{prop:front-jet-substructure}
（测试用合成陈述：内容从略。）
（本命题只涉及部分子态。）
\uses{def:jets-algorithms, prop:pdf-timelike-dglap}
\end{proposition}
\begin{proofsketch}测试。\end{proofsketch}
\paragraph{地位与缺口}测试：见文献。

\begin{proposition}[能量关联子的共线因子化与共形场论中的光线 OPE \tagA]\label{prop:front-lightray-ope}
（测试用合成陈述：内容从略。）
\uses{def:jets-algorithms, prop:pdf-timelike-dglap}
\end{proposition}
\begin{proofsketch}测试。\end{proofsketch}
\paragraph{地位与缺口}测试：见文献。

\section{第 15 节（测试）}\label{sec:front-qed}

\begin{proposition}[NNLO 的红外极点与规范完备的八重态算符 \tagP]\label{prop:front-nrqcd-ir}
（测试用合成陈述：内容从略。）
（本命题只涉及部分子态。）
\uses{def:heavy-nrqcd}
\end{proposition}
\paragraph{地位与缺口}测试：见文献。

\begin{remark}[唯象张力与替代方案]\label{rem:front-quarkonium-alternatives}
（测试用合成陈述：内容从略。）
% uses: none
\end{remark}

\begin{theorem}[由结构函数确定质子的光子 PDF（LUXqed） \tagA]\label{thm:front-luxqed}
（测试用合成陈述：内容从略。）
非微扰地位：\cref{asm:intro-wightman}、\cref{asm:intro-composite}、\cref{asm:intro-bridge}。
\uses{def:pdf-gluon, thm:disall-factorization}
\end{theorem}
\begin{proofsketch}测试。\end{proofsketch}
\paragraph{地位与缺口}测试：本陈述的地位由\cref{thm:disall-factorization}决定。

\begin{example}[点状靶极限]\label{ex:front-luxqed-pointlike}
（测试用合成陈述：内容从略。）
\uses{thm:front-luxqed}
\end{example}

\begin{proposition}[轻子与电子的 PDF、QED 与混合演化核 \tagA]\label{prop:front-lepton-pdf}
（测试用合成陈述：内容从略。）
非微扰地位：\cref{asm:intro-bridge}。
\uses{thm:front-luxqed, thm:pdf-dglap}
\end{proposition}
\begin{proofsketch}测试。\end{proofsketch}
\paragraph{地位与缺口}测试：见文献。

\section{第 16 节（测试）}\label{sec:front-status}

\begin{proposition}[电弱 PDF 与 $Q\gg M_W$ 的电弱对数 \tagA]\label{prop:front-ew-pdf}
（测试用合成陈述：内容从略。）
非微扰地位：\cref{asm:intro-bridge}。
\uses{thm:pdf-dglap}
\end{proposition}
\begin{proofsketch}测试。\end{proofsketch}
\paragraph{地位与缺口}测试：见文献。

\begin{proposition}[电弱 Bloch–Nordsieck 抵消的失效与测量弱荷的观测量 \tagV]\label{prop:front-ew-violation}
（测试用合成陈述：内容从略。）
\uses{def:ir-safe}
\end{proposition}
\paragraph{地位与缺口}测试：本陈述的地位由\cref{prop:ir-bn-failure}决定。

\begin{remark}[开放问题]\label{rem:front-open}
（测试用合成陈述：内容从略。）
% uses: none
\end{remark}
